\section{Discussion}
\label{sec:discussion}

\subsection{Why the band looks the way it does}
\label{sec:discussion-why}

The band of Banco Andino spans the interval from $\bdLowerDf$ to $\bdWorstDf$ of the headline, and the two ends are far from
everything that a plausible dependence produces. The largest capital among the ten planted structures is $\bdTruthMaxOverHeadline$
of the headline, reached by a $t$ copula with 3 degrees of freedom and a crisis matrix, and the worst case is $\bdWorstDf$.
The reason is in the construction of the sharp bounds. The worst case moves the largest losses of each risk type into the
scenarios where the other types have small ones, so that every scenario of the tail has nearly the same total
(Figure~\ref{fig:rearrangement}); no positive dependence structure does that, and it is the structure that makes the value at risk
exceed the sum of the values at risk by $\bdWorstOverComo$ percent. It is a bound that states what cannot be excluded without
knowledge of the dependence, and the planted structures confirm that it is not a scenario. The lower bound is the mirror image, a
structure in which the large losses never coincide, and the independent risks, at $\bdTruthDfGaussianIndependent$ of the
headline, are still far above it.

The interior rows describe the plausible structures, and what they cover is decided by the matrices. With the base matrix alone the
five interior rows contain two of the ten structures; with the stressed matrix added they contain eight, and the interval grows
by $\bdAblWidthStress$ of the capital. The matrix, not the family, is the first-order component, and Proposition~\ref{prop:sqrt}
shows why for the square-root row: the effect of a correlation is the product of the two capital figures, so the pair of the largest
marginals, credit and market, carries most of the movement. The family matters in a second order, through the tail dependence that
Proposition~\ref{prop:tails} describes, and the experiment with the planted structures shows the order of magnitude: a Gaussian
copula calibrated on the exact rank correlation of the structure is within $\bdKGaussMax$ percent for the five Gaussian-family structures and
\bdKTailMin{} to \bdKTailMax{} percent below for the five with tail dependence.

\subsection{Failure cases and negative results}
\label{sec:discussion-failures}

Two predictions were refuted and both are informative. The interior rows of the base matrix covered 2 of the 10 structures
against the 75 percent predicted. The prediction assumed that the interior rows, which are five different aggregation
methods, would bracket the plausible structures; in fact they bracket the plausible structures only when the matrix is right. A
validator that reads the base band alone concludes that the aggregate is between CLP \bdGauss{} and \bdTFour{} billion (the Gaussian
row to the $t$ row of 4 degrees of freedom), and eight of the ten planted structures, which a bank may
reasonably worry about, lie outside that interval. The stress matrix makes the difference. This is a finding about the protocol:
the band must carry the stressed matrix as a first-class component, which the rule of Section~\ref{sec:method-rule} does by naming
the matrices as conventions of the same document.

The second refutation, that more than half of the tests run on the Gaussian row alone, shows that the number of tests is not the
measure of the protection that the band offers. Six tests verify arithmetic, units, consistency and a reference and run on any
single row. The three that read the other rows are the ones that measure the dependence uncertainty, and those fail without the band.
The refutation also gives a use for the table of Section~\ref{sec:ablation-tests}: it says which tests are protections against a
single-figure report.

Three other limits of the results are negative in a milder sense. The additivity of the Euler allocation holds to
$\bdEulerMaxGap$ because it is an identity of the estimator; the test guards the implementation, and the information in the allocation
is its sampling error, which has a standard deviation of $\bdShareSdMax$ percentage points for operational risk. The hypothesis that the sharp range
covers every planted structure was predicted correctly and carries no information about the structures, as
Proposition~\ref{prop:coverage} shows. And the $t$ curve is flat between 2.5 and 3 degrees of freedom within the noise, so the
conjecture of monotonicity is supported on a grid that has one flat step.

\subsection{What the band lets a validator say}
\label{sec:discussion-validator}

The band changes the form of the statement that a validator can make about an aggregate. Without it, the statement is that the
aggregate is a figure from a method with a matrix. With it, the validator can say what share of the width belongs to the choice of
method (the Gaussian and the $t$ rows differ by up to $\bdCopulaDiffFour$ percent, amber in the copula test), what share to
the matrix (the stress moves the Gaussian row by $\bdGaussMove$ percent, amber in the stress test), and what share to the
dependence that no copula of the set spans (the distance from the interior rows to the bounds). The model-risk measures of the band
put numbers on the same: $\mathrm{AM}$ of $\bdAm$ percent says that the worst case is $\bdAm$ percent above the Gaussian figure,
and $\mathrm{RM}$ of $\bdRm$ places it within the band.

The findings also bear on what a bank should report as capital. The case keeps the simple sum as the headline, on the supervisory
position that diversification between risk types is not recognised; the experiments show that the position is conservative for
every planted structure with positive dependence, since the simple sum exceeds the capital of each by $\bdSumErrMin$ to $\bdSumErrMaxOther$ percent, except for
the independent structure, where it exceeds it by $\bdSumMaxErr$ percent. A bank that did claim diversification would have to
choose one row of the band, and the experiments give it a reason for the choice: the Gaussian row underestimates by $\bdGaussDepMin$ to $\bdGaussDepMax$ percent
every planted structure with a stressed matrix or with tail dependence, seven of the ten, and the interior rows of the stressed matrix come closer: on these seven structures the Gaussian row of the stressed matrix falls short by
$\bdGaussStressDepMin$ to $\bdGaussStressDepMax$ percent, against $\bdGaussDepMin$ to $\bdGaussDepMax$ percent below for the Gaussian row of the base matrix.

\subsection{Relation to the published findings}
\label{sec:discussion-literature}

The overestimate of the additive rule over the Gaussian copula in the case, $\bdAdditiveOver$ percent, lies in the range of 30 to 40
percent that Rosenberg and Schuermann report \citep{rosenberg2006}. The same authors find that the choice between a normal and a Student $t$ copula has a more modest effect
on risk than the additive overestimate; the case agrees in order of magnitude, since the largest difference between a $t$ row and
the Gaussian row is $\bdCopulaStat$ percent, against $\bdAdditiveOver$ percent for the simple sum. The excess of the worst value at risk over the comonotonic one,
$\bdWorstOverComo$ percent here, is the phenomenon that Embrechts, Puccetti and R\"uschendorf describe and that makes the simple
sum a non-bound \citep{epr2013}. The dependence of the capital on the degrees of freedom is that of the tail dependence of
Demarta and McNeil \citep{demarta2005}. The Joint Forum's reluctance to rely on aggregation models \citep{jointforum2010} is, in
these numbers, a reluctance to rely on one row: the band of a single case shows that the figures of the rows differ by amounts
that no supervisor can ignore. None of these is a new finding. What the experiments add is the measurement of what the band contains
when the structure is known, and the identification of the matrix as the component that decides it.
